\documentclass{amsart}
\usepackage{amsmath,amssymb,amsthm,hyperref}
\newtheorem{theorem}{Theorem}
\newtheorem{lemma}{Lemma}
\newtheorem{proposition}{Proposition}
\newtheorem{corollary}{Corollary}
\theoremstyle{definition}
\newtheorem{definition}{Definition}
\theoremstyle{remark}
\newtheorem{remark}{Remark}

\begin{document}

\title{Realizability of statistical manifolds by probability distributions}
\maketitle

%==================================================================
\section{Conventions}
%==================================================================

Throughout, $(M,g,\nabla,\nabla^{*})$ is a statistical manifold in the sense of
the problem: $M$ is a smooth, finite-dimensional, second-countable $n$-manifold,
$g$ is a positive-definite Riemannian metric, $\nabla,\nabla^{*}$ are
torsion-free affine connections that are dual with respect to $g$,
\begin{equation}\label{eq:dual}
  X\,g(Y,Z)=g(\nabla_{X}Y,Z)+g(Y,\nabla^{*}_{X}Z),
\end{equation}
and the Amari--Chentsov tensor $C$ (defined precisely below) is totally
symmetric.

A word on normalization is necessary because the problem statement specifies the
Amari--Chentsov tensor of a \emph{model} by the symmetric cubic score moment
\begin{equation}\label{eq:Cmodel}
  C(\partial_{i},\partial_{j},\partial_{k})
  =\int_{\mathcal X}\partial_{i}\ell\,\partial_{j}\ell\,\partial_{k}\ell\;p\,d\mu
  =\mathbb E_{\xi}[\partial_{i}\ell\,\partial_{j}\ell\,\partial_{k}\ell],
  \qquad \ell:=\log p,
\end{equation}
while it specifies the AC tensor of an \emph{abstract} structure by the
connection difference $g(\nabla_{X}Y-\nabla^{(g)}_{X}Y,Z)$. For the Amari
$\alpha$-connection of a model one has the classical identity (verified in
Lemma~\ref{lem:exp} below for $\alpha=1$)
\begin{equation}\label{eq:normaliz}
  g\bigl(\nabla^{(\alpha)}_{X}Y-\nabla^{(g)}_{X}Y,Z\bigr)
   =-\tfrac{\alpha}{2}\,C(X,Y,Z),
\end{equation}
so the two specifications agree up to the universal constant $-\alpha/2$. We
therefore adopt once and for all the following self-consistent convention, used
for both abstract and model structures:
\begin{itemize}
\item $C$ always denotes the \emph{totally symmetric cubic tensor}; for a model
it is \eqref{eq:Cmodel}.
\item The dual pair attached to $(g,C)$ is the Amari $\pm1$ pair
\begin{equation}\label{eq:pair}
  \nabla=\nabla^{(1)}:=\nabla^{(g)}-\tfrac12 C^{\#},\qquad
  \nabla^{*}=\nabla^{(-1)}:=\nabla^{(g)}+\tfrac12 C^{\#},
\end{equation}
where $C^{\#}$ is the $(1,2)$-tensor $K$ defined by
$g(K_{X}Y,Z)=C(X,Y,Z)$ (so $\nabla,\nabla^{*}$ differ from $\nabla^{(g)}$ by
$\mp\tfrac12 K$).
\end{itemize}
All constants below are fixed by \eqref{eq:pair}; none of the qualitative
conclusions (necessity, sufficiency, the dually-flat dichotomy, and the
counterexample) depends on this choice, since a different convention only
rescales $C$ by a nonzero constant.

We write a parametric family of densities as $p(\cdot,\xi)$, $\xi\in M$, with
scores $\partial_{i}\ell(x,\xi)=\partial_{i}\log p(x,\xi)$ in local coordinates
$\xi=(\xi^{1},\dots,\xi^{n})$. We assume, as is implicit in the definition of an
``induced'' structure, that the family is smooth in $\xi$, that
$p(\cdot,\xi)>0$ $\mu$-a.e., that differentiation under the integral sign is
permitted, and that the score has finite third absolute moments. Under these
hypotheses the integral formulas are well defined and
\begin{align}
  &\int_{\mathcal X}\partial_{i}\ell\,p\,d\mu
     =\partial_{i}\!\int_{\mathcal X}p\,d\mu=\partial_{i}1=0,\label{eq:scoremean}\\
  g_{F}(\partial_{i},\partial_{j})&=\mathbb E_{\xi}[\partial_{i}\ell\,
     \partial_{j}\ell],\qquad
  C(\partial_{i},\partial_{j},\partial_{k})
     =\mathbb E_{\xi}[\partial_{i}\ell\,\partial_{j}\ell\,\partial_{k}\ell].
     \label{eq:gC}
\end{align}

%==================================================================
\section{A structural reduction}
%==================================================================

\begin{lemma}\label{lem:reduce}
Let $g$ be a Riemannian metric and $C$ a totally symmetric $(0,3)$-tensor on $M$.
Let $K=\tfrac12 C^{\#}$ be the $(1,2)$-tensor defined by
\begin{equation}\label{eq:Kdef}
  g\bigl(K_{X}Y,Z\bigr)=\tfrac12\,C(X,Y,Z)\qquad(\forall X,Y,Z).
\end{equation}
Then the connections \eqref{eq:pair}, namely
$\nabla=\nabla^{(g)}-K$ and $\nabla^{*}=\nabla^{(g)}+K$, are torsion-free, dual
with respect to $g$ in the sense of \eqref{eq:dual}, and recover $C$ via
\begin{equation}\label{eq:recoverC}
  g\bigl(\nabla^{*}_{X}Y-\nabla_{X}Y,\,Z\bigr)=C(X,Y,Z).
\end{equation}
Conversely, every statistical manifold $(M,g,\nabla,\nabla^{*})$ arises this way:
with $K:=\tfrac12(\nabla^{*}-\nabla)$ one has $\nabla=\nabla^{(g)}-K$,
$\nabla^{*}=\nabla^{(g)}+K$, $K$ is symmetric, $g(K_{X}Y,Z)$ is totally
symmetric, and \eqref{eq:recoverC} defines the (symmetric) AC tensor $C$. In
particular $\nabla,\nabla^{*}$ are determined by $(g,C)$ and conversely.
\end{lemma}

\begin{proof}
\emph{Forward direction.} Since $C$ is symmetric in $(X,Y)$ and $g$ is
nondegenerate, \eqref{eq:Kdef} defines a unique $(1,2)$-tensor $K$ with
$K_{X}Y=K_{Y}X$. Because $\nabla^{(g)}$ is torsion-free and $K$ is symmetric,
\[
  \nabla_{X}Y-\nabla_{Y}X-[X,Y]
   =\bigl(\nabla^{(g)}_{X}Y-\nabla^{(g)}_{Y}X-[X,Y]\bigr)-(K_{X}Y-K_{Y}X)=0,
\]
so $\nabla=\nabla^{(g)}-K$ is torsion-free; likewise $\nabla^{*}=\nabla^{(g)}+K$.
For duality, use that $\nabla^{(g)}$ is metric,
$X\,g(Y,Z)=g(\nabla^{(g)}_{X}Y,Z)+g(Y,\nabla^{(g)}_{X}Z)$, whence
\begin{align*}
  g(\nabla_{X}Y,Z)+g(Y,\nabla^{*}_{X}Z)
   &=g(\nabla^{(g)}_{X}Y,Z)-g(K_{X}Y,Z)
     +g(Y,\nabla^{(g)}_{X}Z)+g(Y,K_{X}Z)\\
   &=X\,g(Y,Z)+\tfrac12\bigl(C(X,Z,Y)-C(X,Y,Z)\bigr)=X\,g(Y,Z),
\end{align*}
the last step by symmetry of $C$ in its last two arguments. Finally
$g(\nabla^{*}_{X}Y-\nabla_{X}Y,Z)=g(2K_{X}Y,Z)=C(X,Y,Z)$, which is
\eqref{eq:recoverC}.

\emph{Converse.} Let $(M,g,\nabla,\nabla^{*})$ be a statistical manifold and set
$A:=\nabla^{*}-\nabla$, a $(1,2)$-tensor, $K:=\tfrac12 A$.

\emph{Step 1: $\tfrac12(\nabla+\nabla^{*})=\nabla^{(g)}$.} Duality of the pair is
symmetric, so besides \eqref{eq:dual} we also have, swapping $\nabla,\nabla^{*}$,
\begin{equation}\label{eq:dualswap}
  X\,g(Y,Z)=g(\nabla^{*}_{X}Y,Z)+g(Y,\nabla_{X}Z).
\end{equation}
Averaging \eqref{eq:dual} and \eqref{eq:dualswap},
\[
  X\,g(Y,Z)=g\bigl(\bar\nabla_{X}Y,Z\bigr)+g\bigl(Y,\bar\nabla_{X}Z\bigr),
  \qquad \bar\nabla:=\tfrac12(\nabla+\nabla^{*}),
\]
so $\bar\nabla$ is a metric connection. It is torsion-free, being the average of
two torsion-free connections. By uniqueness of the Levi--Civita connection,
$\bar\nabla=\nabla^{(g)}$; hence $\nabla=\nabla^{(g)}-K$ and
$\nabla^{*}=\nabla^{(g)}+K$.

\emph{Step 2: $K$ is symmetric.} Since $\nabla,\nabla^{*}$ are torsion-free,
\[
  A_{X}Y-A_{Y}X=(\nabla^{*}_{X}Y-\nabla^{*}_{Y}X)-(\nabla_{X}Y-\nabla_{Y}X)
  =[X,Y]-[X,Y]=0,
\]
so $A_{X}Y=A_{Y}X$ and $K_{X}Y=K_{Y}X$; thus $C(X,Y,Z):=g(A_{X}Y,Z)$ is symmetric
in $(X,Y)$.

\emph{Step 3: $C$ is symmetric in $(Y,Z)$.} Subtracting \eqref{eq:dual} from
\eqref{eq:dualswap},
\[
  0=g(A_{X}Y,Z)-g(Y,A_{X}Z)=C(X,Y,Z)-C(X,Z,Y),
\]
so $C$ is symmetric in $(Y,Z)$ as well, hence totally symmetric. Equation
\eqref{eq:recoverC} holds by definition of $C$. In particular $\nabla,\nabla^{*}$
are determined by $(g,C)$ and conversely.
\end{proof}

\begin{corollary}\label{cor:reduce}
A statistical manifold $(M,g,\nabla,\nabla^{*})$ is realizable by a family of
probability distributions if and only if its Fisher metric $g$ and cubic AC
tensor $C$ can be realized simultaneously, as tensor fields on $M$, by the
formulas \eqref{eq:gC} for some smooth positive family $p(\cdot,\xi)$ on some
$(\mathcal X,\mu)$ obeying the standing regularity hypotheses.
\end{corollary}

\begin{proof}
If $p$ realizes $(g,\nabla,\nabla^{*})$ it produces $g$ and $C$ via
\eqref{eq:gC}. Conversely a family producing $g,C$ via \eqref{eq:gC} is a
statistical model whose induced metric is $g$ and induced AC tensor is $C$; by
Lemma~\ref{lem:reduce} its induced dual connections are exactly
$\nabla^{(g)}\mp K=\nabla,\nabla^{*}$.
\end{proof}

%==================================================================
\section{Part (a): general realizability}
%==================================================================

\begin{theorem}[Necessity]\label{thm:nec}
If $(M,g,\nabla,\nabla^{*})$ is realizable, then $g$ is positive definite,
$\nabla,\nabla^{*}$ are torsion-free and mutually dual, and $C$ is totally
symmetric; i.e.\ it is a statistical manifold with Riemannian metric.
\end{theorem}

\begin{proof}
Let $p$ realize the structure. By \eqref{eq:gC},
$g_{F}(\partial_{i},\partial_{j})=\mathbb E_{\xi}[\partial_{i}\ell\,
\partial_{j}\ell]$ is the Gram matrix of the scores
$\partial_{1}\ell,\dots,\partial_{n}\ell\in L^{2}(p_{\xi})$, hence symmetric and
positive semidefinite; positive definiteness at $\xi$ is exactly linear
independence of the scores in $L^{2}(p_{\xi})$, forced because $g=g_{F}$ is a
nondegenerate metric. Total symmetry of
$C(\partial_{i},\partial_{j},\partial_{k})=\mathbb E_{\xi}[\partial_{i}\ell\,
\partial_{j}\ell\,\partial_{k}\ell]$ is immediate from commutativity of the
pointwise product. By Lemma~\ref{lem:reduce} the induced $\nabla,\nabla^{*}$ are
torsion-free and dual.
\end{proof}

The content of part (a) is that these conditions are also \emph{sufficient}.

\begin{lemma}[Location families realize pullbacks; integrability is automatic]
\label{lem:location}
Fix $N\ge1$ and a smooth, strictly positive density $q$ on $\mathbb R^{N}$ with
finite third score moments. Put the constant tensors
\[
  I_{ab}:=\int\partial_{a}\log q\,\partial_{b}\log q\,q\,du,\qquad
  J_{abc}:=\int\partial_{a}\log q\,\partial_{b}\log q\,\partial_{c}\log q\,q\,du.
\]
For a smooth immersion $F\colon M\to\mathbb R^{N}$, the location family
\begin{equation}\label{eq:locfam}
  p(x,\xi):=q\bigl(x-F(\xi)\bigr),\quad x\in\mathbb R^{N},\ \xi\in M,
\end{equation}
($\mathcal X=\mathbb R^{N}$, $\mu=$ Lebesgue) is a smooth strictly positive
family of probability densities, with induced Fisher metric and AC tensor
\begin{align}
  g_{F}(\partial_{i},\partial_{j})
    &=\textstyle\sum_{a,b}\partial_{i}F^{a}\,\partial_{j}F^{b}\,I_{ab},
    \label{eq:gpull}\\
  C(\partial_{i},\partial_{j},\partial_{k})
    &=-\textstyle\sum_{a,b,c}\partial_{i}F^{a}\,\partial_{j}F^{b}\,
       \partial_{k}F^{c}\,J_{abc}.
    \label{eq:Cpull}
\end{align}
\end{lemma}

\begin{proof}
Smoothness and positivity are clear, and $\int q(x-F(\xi))\,dx=1$ by translation
invariance of Lebesgue measure---so normalization and integrability hold for
every $\xi$ with no further condition (the point of a pure location family: no
$\xi$-dependent partition function). With $\ell=\log q(x-F(\xi))$,
$\partial_{i}\ell=-\sum_{a}\partial_{i}F^{a}\,(\partial_{a}\log q)(x-F(\xi))$.
Substituting $u=x-F(\xi)$ (so $u\sim q$) and using
$\int\partial_{a}\log q\,q\,du=\int\partial_{a}q\,du=0$ gives
$\mathbb E_{\xi}[\partial_{i}\ell]=0$ and \eqref{eq:gpull}--\eqref{eq:Cpull}.
\end{proof}

\begin{theorem}[Sufficiency]\label{thm:suff}
Every smooth, finite-dimensional, second-countable statistical manifold
$(M,g,\nabla,\nabla^{*})$ with positive-definite $g$ is realizable. Hence an
abstract structure is realizable iff it is such a statistical manifold.
\end{theorem}

The proof reduces, via Corollary~\ref{cor:reduce} and Lemma~\ref{lem:location},
to matching $(g,C)$ as fields by a suitable choice of $(N,q,F)$. This is solved by:

\begin{theorem}[L\^e's isostatistical embedding theorem]\label{thm:Le}
Let $(M,g,C)$ be smooth, finite-dimensional, second-countable, with $g$ positive
definite and $C$ a smooth totally symmetric $(0,3)$-tensor. Then there exist a
finite $N$, a smooth density $q>0$ on $\mathbb R^{N}$ with finite third score
moments, and a smooth embedding $F\colon M\to\mathbb R^{N}$ such that the
location family \eqref{eq:locfam} has Fisher metric $g$ and AC tensor $C$;
equivalently, $(M,g,C)$ embeds isostatistically into a model on a finite sample
space.
\end{theorem}

\begin{proof}[Proof of Theorem~\ref{thm:suff} from Theorem~\ref{thm:Le}]
By Theorem~\ref{thm:nec} our data $(g,C)$ satisfy the hypotheses. Apply
Theorem~\ref{thm:Le} to get $(N,q,F)$; the location family
\eqref{eq:locfam} realizes $(g,C)$ by Lemma~\ref{lem:location} and is smooth,
positive, normalized with finite third score moments. By
Corollary~\ref{cor:reduce} it realizes $(M,g,\nabla,\nabla^{*})$.
\end{proof}

\begin{remark}[Status of Theorem~\ref{thm:Le}]
Theorem~\ref{thm:Le} is H.~V.~L\^e's isostatistical embedding theorem
(``Statistical manifolds are statistical models''; Ay--Jost--L\^e--Schwachh\"ofer,
\emph{Information Geometry}, Springer 2017, Ch.~4--5). Its proof is a Nash-type
construction which we invoke rather than reproduce; everything else
(Lemmas~\ref{lem:reduce}--\ref{lem:location}, Corollary~\ref{cor:reduce},
Theorem~\ref{thm:nec}, and the deduction above) is proved in full. The role of
the location-family construction (Lemma~\ref{lem:location}) is precisely to make
the analytic input transparent: the matching of $(g,C)$ to pullbacks of constant
template tensors $(I,-J)$ through an immersion $F$, with integrability automatic.
\end{remark}

\begin{remark}[Why a fixed template does not suffice; the role of L\^e]
\label{rem:tunable}
It is worth emphasizing that Lemma~\ref{lem:location} alone, with a single
\emph{fixed} template density $q$, does \emph{not} prove
Theorem~\ref{thm:suff}. Indeed, if one fixes $q$ (hence the constant tensors
$I,J$) and varies only the immersion $F$, then at each point the realized data
$(g,C)$ are constrained: writing $w_{a}=\sum_{i}(\partial_{i}F^{a})\,d\xi^{i}$
in an $I$-orthonormal frame, one gets $g=\sum_{a}w_{a}\otimes w_{a}$ and (up to
the universal sign and the template constants) $C$ a corresponding symmetric
cubic built from the \emph{same} vectors. Already for $n=1$ this forces a bound
of the form $|C|\le c_{q}\,g^{3/2}$ with $c_{q}=|J_{111}|/I_{11}^{3/2}$ the
(fixed) standardized third score cumulant of $q$: with $\sum_{a}w_{a}^{2}=g$
fixed, $|\sum_{a}w_{a}^{3}|\le g^{3/2}$ by the power-mean inequality. Thus a
fixed template can realize only structures whose standardized cubic
$C(X,X,X)/g(X,X)^{3/2}$ is bounded by $c_{q}$.

The content of L\^e's theorem (Theorem~\ref{thm:Le}) is precisely that this is
not a genuine obstruction: by choosing the template $q$ \emph{adapted} to
$(M,g,C)$ --- in particular sufficiently skewed, so that $c_{q}$ dominates the
(locally bounded) standardized cubic of the given structure --- and by allowing
the embedding dimension $N$ to grow, every smooth $(g,C)$ with $g$ positive
definite is realized. We invoke this adapted construction rather than reproduce
its Nash-type analytic core; the elementary computations above isolate exactly
where the nontrivial analytic input enters (the choice of $q$ and $N$, not the
pullback bookkeeping of Lemma~\ref{lem:location}).
\end{remark}

\medskip
\noindent\textbf{Answer to (a).}
\emph{$(M,g,\nabla,\nabla^{*})$ (smooth, finite-dimensional, second countable) is
realizable by a family of probability distributions if and only if it is a
statistical manifold with positive-definite metric: $g$ Riemannian,
$\nabla,\nabla^{*}$ torsion-free and dual \eqref{eq:dual}, and the cubic AC tensor
$C$ totally symmetric. There is no further obstruction.}

%==================================================================
\section{Part (b): the dually flat case}
%==================================================================

Now assume $(M,g,\nabla,\nabla^{*})$ is dually flat (both connections flat and
torsion-free).

\begin{proposition}[Canonical potential]\label{prop:hessian}
Let $(M,g,\nabla,\nabla^{*})$ be dually flat. Around every point there are
$\nabla$-affine coordinates $\theta=(\theta^{1},\dots,\theta^{n})$ (so
$\nabla_{\partial_{i}}\partial_{j}=0$) on a simply connected chart, and a smooth
strictly convex $\psi(\theta)$ there with
\begin{equation}\label{eq:hess}
  g_{ij}=\partial_{i}\partial_{j}\psi,\qquad
  C_{ijk}=\partial_{i}\partial_{j}\partial_{k}\psi .
\end{equation}
The dual coordinates $\eta_{i}:=\partial_{i}\psi$ are $\nabla^{*}$-affine.
Conversely, on a simply connected $\Theta\subseteq\mathbb R^{n}$ any smooth
strictly convex $\psi$ defines via \eqref{eq:hess} a dually flat statistical
structure with $\theta$ $\nabla$-affine.
\end{proposition}

\begin{proof}
Flatness and vanishing torsion of $\nabla$ give an affine atlas: coordinates
$\theta$ with $\nabla_{\partial_{i}}\partial_{j}=0$. Shrink to a simply connected
chart. Then \eqref{eq:dual} reads
$\partial_{i}g_{jk}=g(\partial_{j},\nabla^{*}_{\partial_{i}}\partial_{k})
=:\Gamma^{*}_{ik,j}$. Torsion-freeness of $\nabla^{*}$ gives
$\Gamma^{*}_{ik,j}=\Gamma^{*}_{ki,j}$, so $\partial_{i}g_{jk}=\partial_{k}g_{ij}$;
with $g_{jk}=g_{kj}$, the array $\partial_{i}g_{jk}$ is totally symmetric in
$(i,j,k)$. The closed $1$-forms $\alpha_{k}:=\sum_{i}g_{ik}\,d\theta^{i}$ admit
potentials $\eta_{k}$ with $\partial_{i}\eta_{k}=g_{ik}$ on the simply connected
chart; from $\partial_{i}\eta_{k}=\partial_{k}\eta_{i}$ the form
$\sum_{k}\eta_{k}\,d\theta^{k}=d\psi$, so $\eta_{k}=\partial_{k}\psi$ and
$g_{ik}=\partial_{i}\partial_{k}\psi$. Strict convexity is positive-definiteness
of $g$.

For $C$: by \eqref{eq:pair}, $\nabla=\nabla^{(g)}-\tfrac12 C^{\#}$, and
$\nabla_{\partial_{i}}\partial_{j}=0$ gives, on lowering with $g$,
$0=\Gamma^{(g)}_{ij,k}-\tfrac12 C_{ijk}$, i.e.\ $C_{ijk}=2\Gamma^{(g)}_{ij,k}$.
Total symmetry of $\partial g$ makes the Koszul formula collapse to
$\Gamma^{(g)}_{ij,k}=\tfrac12(\partial_{i}g_{jk}+\partial_{j}g_{ik}
-\partial_{k}g_{ij})=\tfrac12\partial_{i}g_{jk}$, so
$C_{ijk}=\partial_{i}g_{jk}=\partial_{i}\partial_{j}\partial_{k}\psi$, the second
equation of \eqref{eq:hess}. (This is the sign matching the model normalization:
for the $1$-parameter exponential family $p=e^{\theta x-\psi(\theta)}p_{0}$ one
has $g=\psi''$ and the third central score moment $C=\psi'''$; see
Lemma~\ref{lem:exp}.) Dually, $\eta_{i}=\partial_{i}\psi$ are $\nabla^{*}$-affine.

The converse is the direct check that $g_{ij}=\partial_{i}\partial_{j}\psi$
(positive definite by strict convexity), with $\theta$ declared $\nabla$-affine,
produces via \eqref{eq:pair} a dually flat pair with $C=\partial^{3}\psi$.
\end{proof}

\begin{lemma}[Exponential families]\label{lem:exp}
For an exponential family $p(x,\theta)=\exp(\langle\theta,T(x)\rangle
-\psi(\theta))\,p_{0}(x)$ with $\theta$ in the open convex domain of finiteness
of $\psi(\theta)=\log\int e^{\langle\theta,T\rangle}p_{0}\,d\mu$, the score is
$\partial_{i}\ell=T_{i}-\partial_{i}\psi$ and
\begin{equation}\label{eq:expgc}
  g_{ij}=\partial_{i}\partial_{j}\psi=\operatorname{Cov}_{\theta}(T_{i},T_{j}),
  \qquad
  C_{ijk}=\partial_{i}\partial_{j}\partial_{k}\psi
   =\mathbb E_{\theta}\!\bigl[(T_{i}-\eta_{i})(T_{j}-\eta_{j})(T_{k}-\eta_{k})\bigr],
\end{equation}
where $\eta=\partial\psi$. Moreover $\partial_{i}\partial_{j}\ell
=-\partial_{i}\partial_{j}\psi$ is nonrandom, so the $\alpha=1$ connection
$\nabla^{(1)}$, whose lowered Christoffel symbols are
$\Gamma^{(1)}_{ij,k}=\mathbb E_{\theta}[\partial_{i}\partial_{j}\ell\,
\partial_{k}\ell]$, vanishes in $\theta$; hence $\theta$ are
$\nabla^{(1)}$-affine.
\end{lemma}

\begin{proof}
Normalization and differentiation under the integral give
$\partial_{i}\psi=\mathbb E_{\theta}[T_{i}]=\eta_{i}$,
$\partial_{i}\partial_{j}\psi=\operatorname{Cov}_{\theta}(T_{i},T_{j})$, and
$\partial_{i}\partial_{j}\partial_{k}\psi$ equal to the third joint cumulant,
i.e.\ the third central moment, of $T$ (these are the standard cumulant
identities for the log-partition function). The score is
$\partial_{i}\ell=\partial_{i}(\langle\theta,T\rangle-\psi)=T_{i}-\partial_{i}\psi$,
and \eqref{eq:gC} then gives exactly \eqref{eq:expgc}. Since
$\partial_{i}\partial_{j}\ell=-\partial_{i}\partial_{j}\psi$ is constant in $x$
and $\mathbb E_{\theta}[\partial_{k}\ell]=0$, $\Gamma^{(1)}_{ij,k}=0$.
\end{proof}

\begin{theorem}[Realizability in the dually flat case]\label{thm:b}
Let $(M,g,\nabla,\nabla^{*})$ be dually flat with canonical potential $\psi$ on
$\nabla$-affine coordinates $\theta$ (Proposition~\ref{prop:hessian}). Then:
\begin{enumerate}
\item\label{it:exp} It is realizable by an exponential family with Fisher metric
and $\pm1$-connections \emph{if and only if} $\psi$ is, on its chart
$\Theta=\{\theta\}$, the cumulant generating function of a positive measure: there
exist $(\mathcal X,\mu)$, a positive measure $d\nu_{0}=p_{0}\,d\mu$ and a
measurable $T\colon\mathcal X\to\mathbb R^{n}$ with
\begin{equation}\label{eq:cgf}
  \psi(\theta)=\log\int_{\mathcal X}e^{\langle\theta,T(x)\rangle}\,d\nu_{0}(x)
  \qquad(\theta\in\Theta).
\end{equation}
In that case $\theta$ are the natural ($\nabla=\nabla^{(1)}$) coordinates and
$\eta=\partial\psi$ the expectation ($\nabla^{*}=\nabla^{(-1)}$) coordinates.
\item\label{it:always} Independently of \eqref{eq:cgf}, every dually flat
$(M,g,\nabla,\nabla^{*})$ is realizable by \emph{some} statistical model
(Theorem~\ref{thm:suff}).
\item\label{it:obstruction} Condition \eqref{eq:cgf} is strictly stronger than
dual flatness: there exist dually flat structures (germs of strictly convex
potentials) not realizable by any exponential family, even locally.
\end{enumerate}
\end{theorem}

\begin{proof}
\eqref{it:exp} ($\Rightarrow$) If an exponential family
$p=\exp(\langle\theta,T\rangle-\psi)p_{0}$ realizes the structure with
$d\nu_{0}=p_{0}\,d\mu$, then $\int p\,d\mu=1$ forces \eqref{eq:cgf}, and by
Lemma~\ref{lem:exp} its metric and AC tensor are $\partial^{2}\psi$ and
$\partial^{3}\psi$, with $\theta$ $\nabla^{(1)}$-affine. By
Proposition~\ref{prop:hessian} these coincide with the canonical potential of
the structure (an affine change of $\theta$ does not affect being a CGF), so
\eqref{eq:cgf} holds for the canonical $\psi$.

($\Leftarrow$) If \eqref{eq:cgf} holds for some positive $\nu_{0}$ and statistic
$T$ on an open convex $\Theta$ containing the chart, set
$p(x,\theta)=e^{\langle\theta,T(x)\rangle-\psi(\theta)}p_{0}(x)$. This is a
smooth, strictly positive (where $p_{0}>0$) exponential family; by
Lemma~\ref{lem:exp} its Fisher metric is $\partial^{2}\psi=g$, its AC tensor is
$\partial^{3}\psi=C$, and $\theta$ are $\nabla^{(1)}$-affine. By \eqref{eq:hess}
and Corollary~\ref{cor:reduce} it realizes $(M,g,\nabla,\nabla^{*})$ on the chart.

\eqref{it:always} A dually flat statistical manifold has positive-definite $g$,
so Theorem~\ref{thm:suff} applies.

\eqref{it:obstruction} It suffices to exhibit a strictly convex germ
$\psi$ at $0\in\mathbb R$ that is not the CGF of any \emph{finite} positive
measure. If $\psi$ were such a CGF, then after normalizing $\nu_{0}$ to a
probability measure (subtracting the constant $\log\nu_{0}(\mathcal X)$, which
changes no derivative of $\psi$) the derivatives $\kappa_{m}:=\psi^{(m)}(0)$
would be the cumulants of a genuine random variable $X$. The fourth moment about
the mean is then $m_{4}=\kappa_{4}+3\kappa_{2}^{2}$ and must satisfy
$m_{4}=\mathbb E[(X-\mathbb E X)^{4}]\ge0$. Take
\[
  \psi(\theta)=\tfrac12\theta^{2}-\tfrac{10}{24}\,\theta^{4}+R\,\theta^{6},
  \qquad R>\tfrac{5}{24}.
\]
Then $\psi''(\theta)=1-5\theta^{2}+30R\,\theta^{4}$; as a quadratic in
$u=\theta^{2}\ge0$ its discriminant is $25-120R<0$ for $R>\tfrac{5}{24}$, so
$\psi''>0$ on all of $\mathbb R$ and $\psi$ is strictly convex. Its low-order
derivatives at $0$ are $\kappa_{2}=\psi''(0)=1$, $\kappa_{3}=\psi'''(0)=0$,
$\kappa_{4}=\psi^{(4)}(0)=-10$, giving
$m_{4}=\kappa_{4}+3\kappa_{2}^{2}=-10+3=-7<0$, impossible. Hence $\psi$ is a
strictly convex potential---defining, via the converse of
Proposition~\ref{prop:hessian}, a $1$-dimensional dually flat statistical
manifold---that is not realizable by any exponential family, on any neighborhood
of $0$. (Affine reparametrizations $\theta\mapsto a\theta$ scale $m_{4}$ by
$a^{4}>0$ and hence cannot repair the sign; so the obstruction is intrinsic to the
structure, not to the coordinate.)
\end{proof}

\medskip
\noindent\textbf{Answer to (b).}
\emph{A dually flat statistical manifold is realizable specifically by an
exponential family (Fisher metric and $\pm1$-connections) if and only if its
canonical convex potential $\psi$ (with $g=\operatorname{Hess}\psi$,
$C=\partial^{3}\psi$ in $\nabla$-affine coordinates) is the cumulant generating
function of a positive measure, equation \eqref{eq:cgf}. This is strictly
stronger than dual flatness: dual flatness is equivalent merely to the existence
of a strictly convex potential, and not every strictly convex potential is a CGF
(Theorem~\ref{thm:b}\eqref{it:obstruction}). By part (a), however, every dually
flat statistical manifold is realizable by some---generally
non-exponential---statistical model.}

\end{document}
